%% notes-tex/019-simb-multimodal/sections/6-carotene.tex
%% [[experiments.019-simb-multimodal.scripts.pigment_noise_ceiling]]
%% SOURCE: ceiling from experiments/019-simb-multimodal/results/pigment_noise_ceiling.json;
%% scores from experiments/019-simb-multimodal/results/round_leaderboards.csv.

\section{Beta-carotene production\secstatus{todo}}
\label{sec:carotene}

\subsection{Why this target is different from the other pigment\secstatus{todo}}

Ozaydin 2013 (doi:10.1016/j.ymben.2012.07.010) screened the deletion collection carrying a
beta-carotene pathway and scored each colony's color \textbf{by eye on an ordinal scale
from $-5$ to $+5$}. The dataset holds 4,474 records, of which 4,344 have a single replicate
and 130 have two or three. That is a different kind of measurement from the Cachera
fluorescence readout, and it needs a different ceiling: a Pearson ceiling on a subjective
ordinal is the wrong object, so rank agreement is used.

Two estimates exist, and they disagree by a factor of seven.

\begin{table}[htbp]
\centering\small
%% SOURCE: experiments/019-simb-multimodal/results/pigment_noise_ceiling.json, beta_carotene
\begin{tabular}{llrr}
\toprule
estimate & basis & $n$ & Spearman \\
\midrule
replicate max against min & 130 replicated rows & 130 & \textbf{0.544} \\
independent re-screen & \file{2ndRoundOfTransformations} sheet & 119 & 0.075 \\
\quad range-restriction corrected & Thorndike case 2 & 119 & 0.105 \\
\bottomrule
\end{tabular}
\caption[]{Reliability of the beta-carotene score. Neither estimate is clean. The
max-against-min comparison is biased in both directions at once: the two are drawn from one
replicate set, which inflates agreement, and they are the widest possible split of that
set, which deflates it. The re-screen is a genuine test-retest but was run on selected top
hits, so its 1st-screen scores sit almost entirely in $3\ldots5$ of a $-5\ldots+5$ scale.}
\label{tab:carotene-ceiling}
\end{table}

\notesec{Why the $0.075$ is not evidence that the score is worthless} The re-screen is the
better \emph{design} of the two, a genuine independent test-retest rather than two draws
from one replicate set. But it was run only on \textbf{selected top hits}, so its first-round
scores sit almost entirely in $3 \ldots 5$ of a $-5 \ldots +5$ scale, with a standard
deviation of $0.89$ against $1.47$ over the full screen (\cref{fig:carotene-ceiling}b).

A correlation computed on a narrowed slice of a population is attenuated by the narrowing
alone, independently of how good the measurement is. The intuition: if every strain
re-screened is already known to be a high producer, the only variation left to correlate is
the variation \emph{among} high producers, which is small and mostly noise. Thorndike's case
2 formula corrects for exactly this, given the restricted and unrestricted standard
deviations, and it takes the $0.075$ to $0.105$. That is still low, but it is a correction
for one restriction, and the sampling is selected on more than range alone.

The reported ceiling is therefore $0.544$, from the replicate comparison, with the caveat
that neither estimate is clean and they disagree by a factor of five even after correction.

\tcfigfit{figures/beta_carotene_noise_ceiling.pdf}{Reliability of the beta-carotene score,
which is rank agreement rather than a Pearson because the target is a hand-scored ordinal.
(a) The primary estimate, comparing the maximum against the minimum score among the 130
strains with more than one replicate, over unrestricted strains: Spearman $0.544$. (b) The
independent re-screen, which is a genuine test-retest but was run only on selected top hits;
the shaded band marks the middle $90\,\%$ of its first-screen scores, which spans a standard
deviation of $0.89$ against $1.47$ over the full screen. That narrowing is the range
restriction responsible for its raw $0.075$, and it is visible here rather than
asserted.}{fig:carotene-ceiling}

\subsection{What was measured\secstatus{todo}}

The strand is 136 runs across three projects. The best validation
\file{spearman_per_feature} is $\mathbf{0.1371}$, run \texttt{5cjpn6zj} in
\file{beta_carotene_v3}, peaking at \textbf{epoch 9 of 49}. The median over all 136 is
$0.0513$. Against the $0.544$ ceiling that is $25\,\%$ realized.

\notesec{Correction, and it is a correction to the method rather than to a judgment call}
Revision 2 reported the best as $0.118$ from run \texttt{o0aadmot} peaking at epoch 799 of
999. That run is only the \emph{third} best in the strand. The leaderboard had fetched
validation curves for a candidate subset of each project, chosen as the top runs by final
logged value and by epoch count, and \texttt{5cjpn6zj} peaks at epoch 9, so it ranks low on
\emph{both} keys at once and was invisible to the rule. \textbf{This is a structural blind
spot, not bad luck}: a run that peaks very early and then decays is exactly what the
candidate rule cannot see, and it falsifies the previous assumption that incomplete run
coverage risked only medians and never maxima (\cref{sec:coverage}).

\begin{keybox}
\textbf{Beta-carotene is limited by its measurement, not by the model.} A hand-scored
ordinal with one replicate per strain has a reliability no architecture recovers, and the
$0.1371$ best \file{spearman_per_feature} sits at $25\,\%$ of a $0.544$ ceiling that is
itself uncertain by a factor of five. The strand carries real signal and is worth keeping
for the coverage argument in \cref{sec:betaxanthin}, where the claim is about which genes a
flux model can represent and does not depend on the absolute score.
\end{keybox}

\notesec{Correction to revision 2 on convergence, which the coverage fix reverses} Revision 2
said the best run ``peaks late (epoch 791 of 999), so it shares the expression strand's
under-training rather than the metabolite strands' early overfitting''. With the true best
run in hand, the opposite is the case: it peaks at \textbf{epoch 9 of 49} and decays. So
beta-carotene groups with the \emph{metabolite} strands, which peak early and then overfit,
not with expression. That is also what makes it a low-value target for more compute: extra
epochs are not what it is short of.

\notesec{Cannot say} Which of the two ceiling estimates is right, and therefore whether
$0.1371$ is a quarter of achievable or most of it. That gap is resolvable only with a
replicate design the source data does not contain.

\notesec{Recommendation} Keep beta-carotene as a second production target for the coverage
argument in \cref{sec:betaxanthin}, where the claim is about which genes a flux model can
represent and does not depend on the absolute score. Do not spend a sweep on lifting its
number.
